\section{Companion zeros and local fields}\label{sec:local-tools}

Throughout this section $b$ is a simple actual zero, $S=p''(b)/p'(b)\ne0$,
and $H,Q,d,D,q,\mu,\sigma$ have the meanings in
\eqref{pr:field-notation} and \eqref{pr:normalized-moments}. Put
\[
 \alpha_j=\frac1{S(b-\xi_j)},\qquad
 J(w)=\frac{S p(b+w/S)}{p'(b)},\qquad K=p''/p'.
\]
Thus $\sum\alpha_j=1$, $\sum|\alpha_j|^2=q$ and
$\max|\alpha_j|=\mu$. In particular $q>0$. The exact identities are
\begin{equation}\label{lt:finite-product}
 J'(w)=\prod_j(1+\alpha_jw),\qquad
 J(w)=\int_0^w\prod_j(1+\alpha_jv)\,dv.
\end{equation}
When a phase or curvature conclusion is asserted, $b$ is assumed
rightmost among all original zeros. The constructed companions are
zeros of $p$.

\subsection{Logarithmic bounds and circles}

For $m\ge2$ and $\mu|w|<1$, the logarithm vanishing at zero satisfies
\begin{equation}\label{lt:log-tail}
 \log J'(w)=\sum_{k=1}^{m-1}\frac{(-1)^{k+1}w^k}{k}
                    \sum_j\alpha_j^k+E_m(w),\qquad
 |E_m(w)|\le\frac{q\mu^{m-2}|w|^m}{m(1-\mu|w|)}.
\end{equation}
Indeed $\sum_j|\alpha_j|^k\le q\mu^{k-2}$ for $k\ge2$, so the
entire remaining logarithmic series is bounded by a geometric series.
Also, throughout $|w|\le Z<1/\mu$,
\begin{equation}\label{lt:log-derivative}
 \left|\frac{J''}{J'}(w)\right|
 \le1+qZ+\frac{q\mu Z^2}{1-\mu Z}.
\end{equation}
This follows from
$J''/J'=1-w\sum\alpha_j^2+w^2\sum\alpha_j^3/(1+\alpha_jw)$.

\begin{lemma}\label{lt:circle}
Suppose $J'$ is nonzero on $|w|\le Z_{\mathrm{out}}$,
$|J''/J'|\le L$ there, and $|\zeta|+R<Z_{\mathrm{out}}$. If
\[
 \frac{|J(\zeta)|}{|J'(\zeta)|}
       +\frac{e^{LR}-1-LR}{L}<R,
\]
then $|w-\zeta|<R$ contains exactly one zero of $J$, which is simple.
\end{lemma}
\begin{proof}
Since $J''(0)/J'(0)=\sum_j\alpha_j=1$, the hypothesis gives $L\ge1$.
On every segment $\zeta+tv$, $|v|\le R$, integration of the
logarithmic derivative gives
$|J'(\zeta+tv)/J'(\zeta)-1|\le e^{Lt|v|}-1$.
A further integration gives
\[
 |J(\zeta+v)-J(\zeta)-vJ'(\zeta)|
 \le |J'(\zeta)|\frac{e^{L|v|}-1-L|v|}{L}.
\]
On $|v|=R$, Rouch\'e's theorem compares
$J(\zeta+v)$ with $vJ'(\zeta)$. The critical-free containing disk
makes its unique zero simple.
\end{proof}

\subsection{Gaussian quartic centers}

Put $\theta=-\sum\alpha_j^2=K'(b)/S^2$, so $|\theta|\le q$.
For $s=\pm2\pi i$ define
\begin{align}
 \zeta_s={}&s+(s-s^2/2)\theta+(3s-2s^2+s^3/2)\theta^2\notag\\
 &+(15s-21s^2/2+23s^3/6-5s^4/8)\theta^3\notag\\
 &+(105s-72s^2+88s^3/3-22s^4/3+7s^5/8)\theta^4.
 \label{lt:gaussian-center}
\end{align}
If $\zeta_s=\sum_{j=0}^4c_j(s)\theta^j$, write
\[
 t\zeta_s+\theta t^2\zeta_s^2/2
       =st+\sum_{j=1}^9 A_j(s,t)\theta^j.
\]
Thus $A_j=tc_j+(t^2/2)\sum_{k+l=j-1}c_kc_l$, where a coefficient
outside $0\le j\le4$ is zero. Define
\[
 E_0=1,\qquad E_j=\frac1j\sum_{k=1}^{\min(j,9)}kA_kE_{j-k},
 \qquad I_0=0,\quad I_m=\frac{1-mI_{m-1}}s\quad(m\ge1).
\]
For a polynomial $f(t)=\sum f_mt^m$ set
$\mathcal I_s f=\sum f_m I_m$. Since $e^s=1$,
$\mathcal I_s f=\int_0^1e^{st}f(t)\,dt$ at the two indicated values
of $s$. The coefficients of the Gaussian primitive at the center are
\[
 a_j(s)=\mathcal I_s\left(\sum_{k=0}^{\min(4,j)}c_kE_{j-k}\right).
\]
The recurrence gives $a_0=\cdots=a_4=0$ exactly. It can alternatively
be used to determine $c_1,\ldots,c_4$ successively: the new center
coefficient enters the corresponding residual with coefficient one.

The following strict bounds hold for
$6\le p\le44/7$, $0\le t\le1$, with $s=\pm ip$:
\begin{center}\small
\begin{tabular}{cl}
\toprule
polynomials & bounds in increasing order\\\midrule
$c_1,\ldots,c_4$ & $21,133,1030,9000$\\
$A_1,\ldots,A_9$ & $63/10,45,400,3400,86000,320000,$\\
 & $1700000,9100000,39000000$\\
$a_5,\ldots,a_{11}$ & $82000,350000,1900000,13000000,$\\
 & $120000000,3100000000,23000000000$\\\bottomrule
\end{tabular}
\end{center}
Here and below a bound for a complex polynomial means a bound for its
modulus. To verify these inequalities, substitute $s=ip$ in the
recurrences and subtract
the squared modulus from the square of the proposed bound. Convert
the resulting rational polynomial to tensor Bernstein form on
$[6,44/7]\times[0,1]$, omitting the second variable when absent.
If every control coefficient is positive, accept the box; otherwise
bisect it and use the exact de Casteljau averages. At depth $d$ the
path-variable axis is chosen unless $d\equiv2\pmod3$, when the
$p$-axis is chosen. For one-variable polynomials bisect the $p$-axis.
All twenty inequalities terminate with positive coefficients on a total
of twenty-four closed boxes, at depth at most twenty. The recurrences
and bounds specify the certificate;
\cert{check_gaussian_domains_new.py} implements it with rational arithmetic.
The coefficient identities $a_0=\cdots=a_4=0$ are checked by exact
polynomial arithmetic. They apply to the primitive at $s=\pm2\pi i$,
where $e^s=1$. The resulting residual polynomials are bounded on
$6\le p\le44/7$.

Set
\begin{align*}
 Z(q)&=63/10+21q+133q^2+1030q^3+9000q^4,\\
 M(q)&=(63/10)q+45q^2+400q^3+3400q^4+86000q^5\\
     &\quad+320000q^6+1700000q^7+9100000q^8+39000000q^9.
\end{align*}
Let $P_{11}$ be the degree-eleven Taylor polynomial of $Z(q)e^{M(q)}$,
and let $B_5,\ldots,B_{11}$ denote the seven residual bounds in the
table. For $J_\theta(w)=\int_0^we^{z+\theta z^2/2}\,dz$ we have
\begin{equation}\label{lt:gaussian-remainder}
 |J_\theta(\zeta_s)|
 <\sum_{j=5}^{11}B_jq^j+Z(q)e^{M(q)}-P_{11}(q).
\end{equation}
Every coefficient in the last difference is nonnegative. Expand the
exponential along the path, bound the retained coefficients by the table
and every higher coefficient by $Z e^M$.
For exact numerical bounds we use
\begin{equation}\label{lt:exp-upper}
 \mathcal E_{24}(x)=\sum_{k=0}^{24}\frac{x^k}{k!}
       +\frac{x^{25}}{25!(1-x/26)}>e^x\qquad(0<x<26).
\end{equation}
The quotient of the right side of \eqref{lt:gaussian-remainder} by
$q^5$ increases with $q$. Substituting \eqref{lt:exp-upper} proves
the respective coefficients $132200,126000,86200,89000$ at
$q=1/20,97/2000,7/625,9/556$.

\begin{proposition}\label{lt:quartic}\label{lt:fine200}\label{lt:wideD100}
Each row of the following table gives two distinct simple actual zeros,
other than $b$, in the disks $|S(z-b)-\zeta_\pm|<R$.
The disks lie in the indicated critical-free containing disk.
\begin{center}\small
\begin{tabular}{ccclc}
\toprule
row&$q_0$&$\mu_0$&$R(q,\mu)$&$Z_{\mathrm{out}}$\\\midrule
coarse&$97/2000$&$1/270$&$300000q^5+1300q\mu$&$41/5$\\
fine&$7/625$&$1/200$&$120000q^5+220q\mu$&$659/100$\\
wide&$9/400$&$1/100$&$250000q^5+400q\mu$&$174/25$\\
origin&$33/1000$&$8/625$&$3/10$&$15/2$\\
medium&$9/556$&$1/139$&$120000q^5+225q\mu$&$168/25$\\\bottomrule
\end{tabular}
\end{center}
Every row holds on $0<q\le q_0$, $0<\mu\le\mu_0$.
At a rightmost zero the fine and medium rows give
$|\Im S|<4Q/H$, and the wide row gives $|\Im S|<(23/5)Q/H$.
\end{proposition}
\begin{proof}
For the first four rows use $J'=e^{w+\theta w^2/2+E_3(w)}$ and
\eqref{lt:log-tail}. The following constants are exact:
\begin{center}\small
\begin{tabular}{c|rrrrrrr}
\toprule
row&$Z_0$&$M_0$&$E_0$&$B$&$P_0$&$L$&$\beta$\\\midrule
coarse&$39/5$&$101/200$&$43/25$&$126000$&$550$&$143/100$&$27/100$\\
fine&$657/100$&$2/25$&$109/100$&$86200$&$175$&$11/10$&$1/100$\\
wide&$343/50$&$43/250$&$5/4$&$132200$&$250$&$59/50$&$3/50$\\
origin&$719/100$&$7/25$&$71/50$&$132200$&$350$&$32/25$&$11/50$\\\bottomrule
\end{tabular}
\end{center}
At each upper corner, set
$r_3=q_0\mu_0 Z_0^3/[3(1-\mu_0Z_0)]$. The finite predicates are
\begin{gather}
 Z(q_0)<Z_0,\quad M(q_0)<M_0,\quad
 \mathcal E_{24}(M_0+r_3)<E_0,\quad
 \frac{Z_0^4E_0}{12(1-\mu_0Z_0)}<P_0,\label{lt:gaussian-gates1}\\
 Z_0+R_0<Z_{\mathrm{out}}<1/\mu_0,\quad
 1+q_0Z_{\mathrm{out}}+
 \frac{q_0\mu_0Z_{\mathrm{out}}^2}{1-\mu_0Z_{\mathrm{out}}}<L,\notag\\
 \frac{\mathcal E_{24}(LR_0)-1-LR_0}{LR_0}<\beta,
 \label{lt:gaussian-gates2}
\end{gather}
where $R_0=R(q_0,\mu_0)$. The center residual obeys
\[
 |J(\zeta_s)/J'(\zeta_s)|<E_0Bq^5+E_0P_0q\mu.
\]
Indeed the primitive discrepancy is bounded by integrating
$q\mu Z_0^3t^3/[3(1-\mu_0Z_0)]$ for $0\le t\le1$; its integral
contains the factor $1/12$, and $|J'(\zeta_s)|>1/E_0$.
For the first three rows the two coefficient inequalities
$E_0B<(1-\beta)[q^5]R$ and
$E_0P_0<(1-\beta)[q\mu]R$ prove the circle criterion; the brackets
denote coefficients of the indicated monomials in $R$.
In the origin
row their sum at the corner is less than $0.217280<(1-\beta)R_0$.
Every majorant is increasing, so these predicates hold on the
closed rectangle. Source exclusion and separation follow respectively
from
\[
 6-21q_0-133q_0^2-1030q_0^3-9000q_0^4>R_0,
\]
\[
 12(1-q_0-17q_0^2-137q_0^3-313q_0^4)>2R_0.
\]
All displayed predicates are rational after \eqref{lt:exp-upper} is
substituted, and direct evaluation gives positive differences.

For the medium row estimate the discrepancy relative to
the endpoint derivative. Write
$g(t)=t\zeta_s+\theta t^2\zeta_s^2/2$. The coefficient of $\theta$
in $g(t)-g(1)-(t-1)s$ is
$(t-1)s+(t^2-t)s^2/2$ and has modulus less than $63/10$.
To see this, put $u=1-t$ and use $p^2=(2\pi)^2<40$ and
\[
 1-u^2[1+10(1-u)^2]=(1-u)(10u^3-10u^2+u+1).
\]
The cubic on the right has Bernstein coefficients
$(1,7/6,1/2,1/4)$ on $[0,1/2]$,
$(1/4,1/8,5/48,11/32)$ on $[1/2,3/4]$, and
$(11/32,7/12,13/12,2)$ on $[3/4,1]$.
For the remaining coefficients use
$(t-1)c_k+(t^2-1)\sum_{i+j=k-1}c_ic_j/2$.
This gives the relative-path majorant
\begin{align*}
 M_r(q)={}&(63/10)q+(2653/10)q^2+(10442/5)q^3+18282q^4\\
 &+(174349/2)q^5+325990q^6+1727450q^7
       +9270000q^8+40500000q^9.
\end{align*}
At $q_0=9/556$, $\mu_0=1/139$, take
\begin{gather*}
 Z_0=6681/1000,\quad M_0=29/250,\quad M_{r,0}=183/1000,\\
 E_0=57/50,\quad L=28/25,\quad\beta=3/200.
\end{gather*}
The bounds $Z(q_0)<Z_0$, $M(q_0)<M_0$, $M_r(q_0)<M_{r,0}$ and
$e^{M_0+r_3}<E_0$ follow from the same rational predicates.
Dividing the integral discrepancy by the actual $J'(\zeta_s)$ gives
\begin{equation}\label{lt:relative-gaussian-error}
 \left|\frac{J(\zeta_s)-J_\theta(\zeta_s)}{J'(\zeta_s)}\right|
 <\frac{q\mu Z_0^4e^{M_{r,0}+2r_3}}{12(1-\mu_0Z_0)}
 <215q\mu.
\end{equation}
The path remainder and endpoint denominator each contribute a factor
$e^{r_3}$. The model contribution is less than
$101460q^5$. Now $101460<(1-\beta)120000$ and
$215<(1-\beta)225$. The other circle, separation and containment
predicates in \eqref{lt:gaussian-gates2} hold with
$Z_{\mathrm{out}}=168/25$.
These estimates prove the medium row.

At a rightmost $b$, both companions satisfy $\Re(w/\sigma)\le0$.
Choose the sign of $s$ for which
$\Re(s/\sigma)=2\pi|\Im\sigma|$. The triangle inequality gives
\[
 2\pi|\Im\sigma|<q\left(21+133q+1030q^2+9000q^3+R/q\right).
\]
At the fine corner the bracket is less than $24$; at the wide corner
it is less than $28.681<(157/25)(23/5)$; at the medium corner it is
less than $4(157/25)$. Since $2\pi>157/25>6$, the phase bounds follow.
\end{proof}

\subsection{Signed curvature from the two companions}

Write $p_0=2\pi$ for this subsection only, and define
\begin{align*}
 B(\theta)&=1+\theta+(3-p_0^2/2)\theta^2
 +(15-23p_0^2/6)\theta^3
 +(105-88p_0^2/3+7p_0^4/8)\theta^4,\\
 T(\theta)&=\theta+4\theta^2+(21-5p_0^2/4)\theta^3
 +(144-44p_0^2/3)\theta^4.
\end{align*}
Then $\zeta_\pm=(p_0^2/2)T\pm ip_0B$. If the actual roots lie
in disks of radius $R$ about these centers, rightmostness gives
\begin{equation}\label{lt:signed-projection}
 \frac{p_0^2}2\Re(T/\sigma)+p_0|\Im(B/\sigma)|<R.
\end{equation}
For other centers, add to $R$ an upper bound for their displacements
from \eqref{lt:gaussian-center}.

\begin{lemma}\label{lt:signed-moment}
For $r=|\theta|\le1/20$, put
\[
 v=r+17r^2+137r^3+313r^4,\quad m=1-v.
\]
\[
 A_r=\frac{1+17r+137r^2+313r^3}{m},\quad
 B_r=\frac{18+154r+450r^2+313r^3}{m}.
\]
\[
 \mathcal C(r)=1+(2+100/199)rA_r^2+3rB_r
        +(600/199)r^2A_r^2+4rA_r+29r+436r^2.
\]
The two companion disks imply
\begin{equation}\label{lt:curvature-Cq}
 \frac{\Re K'(b)}{H^2}<\mathcal C(q)q^2+\frac{2R}{p_0^2}
 \qquad(q\le1/20).
\end{equation}
It also implies the sharper bound with $\mathcal C(q)$ replaced by
$6/5$. For $q\le3/50$ the replacement $13/10$ is valid provided
an independent circle estimate gives
$|\Im\sigma|<8q$.
\end{lemma}
\begin{proof}
Set $\tau=B/|B|=c+is$ and $\delta=\Im(B/\sigma)$.
The coefficient bounds give $|B|\ge m$, $c>99/100$,
$1-c\le(100/199)s^2$, $|s|\le rA_r$, and
$|s-\Im\theta|\le r^2B_r$. For
$f(\sigma)=\Re(\theta\sigma^2-T/\sigma)$,
\[
 |f(\sigma)-f(\tau)|
 \le r(3+4r+29r^2+436r^3)|\sigma-\tau|.
\]
The exact chord identity is
\[
 \frac{|\sigma-\tau|^2}{\delta^2}
   =\frac{2}{|B|^2(1+\Re(\sigma\overline\tau))}.
\]
On the right semicircle, $\Re(\sigma\overline\tau)\ge-|s|$.
The endpoint comparison at $r=1/20$ gives
$|\sigma-\tau|\le(9/5)|\delta|$ and
\[
 (9/5)r(3+4r+29r^2+436r^3)<7/22<2/p_0.
\]
The case $\delta=0$ follows by continuity: the antipodal point
$-\tau$ has negative real part.

Writing $\theta=x+iy$, expansion of $f(\tau)$ gives
\begin{align*}
 y^2-4x^2&+x(1-c)-2xs^2-y[(s-y)(2c+1)+2y(c-1)]\\
 &+4(x^2-y^2)(1-c)-8xys
 -(21-5p_0^2/4)\Re(\theta^3\overline\tau)
 -(144-44p_0^2/3)\Re(\theta^4\overline\tau).
\end{align*}
Use the bounds above, $|xy|\le r^2/2$, and the coefficient bounds
$|21-5p_0^2/4|<29$, $|144-44p_0^2/3|<436$.
The result is $f(\tau)\le\mathcal C(r)r^2$.
Every term in $\mathcal C$ increases, and
$\mathcal C(97/2000)<9$, $\mathcal C(7/625)<11/5$.
Combining this estimate with \eqref{lt:signed-projection} cancels
$\delta$ and proves \eqref{lt:curvature-Cq}.

For the two sharper constants there is a finite polynomial proof.
Write $\theta=re^{i\phi}$, $u=\cos\phi$, and let $T_j(u)$ be the
Chebyshev polynomials. If $B=\sum b_j\theta^j$ and
$T=\sum d_j\theta^j$, define
\begin{align*}
 M&=\sum_{i,j}b_ib_jr^{i+j}T_{|i-j|}(u),\\
 U&=\sum_{i,j}b_ib_jr^{i+j+1}T_{i+j+1}(u),\qquad
 V=\sum_{i,j}d_ib_jr^{i+j}T_{|i-j|}(u),\\
 P&=(U-V)/r^2,\qquad W=V(M-1)/r^2.
\end{align*}
The last two expressions are polynomials, including at $r=0$.
For $(r_0,k)=(1/20,6/5)$ or $(3/50,13/10)$, let
$m_0=1-r_0-17r_0^2-137r_0^3-313r_0^4$. The polynomials
\begin{equation}\label{lt:signed-polynomials}
 A=kM-P,\qquad (1+m_0)A+W
\end{equation}
are positive on $0\le r\le r_0$, $-1\le u\le1$ and
\[
 3947841760435743/10^{14}\le p_0^2
           \le3947841760435744/10^{14}.
\]
Indeed, tensor Bernstein conversion on the three closed $u$ intervals
$[-1,0]$, $[0,1/2]$, $[1/2,1]$ gives 4950 positive rational
coefficients for each $r_0$. The conversion is exactly
$x^k=\sum_{j=k}^n\binom jk\binom nk^{-1}
\binom njx^j(1-x)^{n-j}$ after affine substitution in each variable.
These polynomials and the three boxes specify the certificate.
Machin's formula and alternating arctangent sums of
lengths twenty and twenty-one give the stated enclosure of $p_0^2$.
Now
\[
 f(\tau)/r^2=\frac1M\left(P-\frac{W}{1+\sqrt M}\right).
\]
If $W\ge0$ use $A>0$; if $W<0$ use $1+\sqrt M\ge1+m_0$ and
the second inequality in \eqref{lt:signed-polynomials}. Thus
$f(\tau)\le kr^2$ on the whole disk, with $r=0$ supplied by continuity.

At the larger cap $r\le q\le3/50$, the phase bound gives
$\Re\sigma>7/8$. The coefficient bounds give
$|B|>21/25$, $|\Im\tau|<1/5$, $\Re\tau>97/100$, hence
$\Re(\sigma\overline\tau)>3011/4000$.
The rational comparison
$(21/25)^2(1+3011/4000)(4/3)^2>2$ gives
$|\sigma-\tau|<(4/3)|\delta|$. At $r=3/50$,
\[
 (4/3)r(3+4r+29r^2+436r^3)=214911/781250<7/22<2/p_0.
\]
The same cancellation proves the second assertion on this larger domain.
\end{proof}

Two consequences needed by the field argument in Section~\ref{sec:matrix-fields}
are
\begin{align}
 \Re K'(b)/H^2&<\mathcal C(q)q^2+(50000/3)q^5+(650/9)q\mu
 &&\text{on the coarse row},\label{lt:coarse-curvature}\\
 \Re K'(b)/H^2&<(9/4)q^2+13q\mu
 &&\text{on the fine row}.\label{lt:fine-curvature}
\end{align}
These follow from the companion disks, since $p_0>6$ and
$11/5+(120000/18)(7/625)^3<9/4$, $220/18<13$.
On the medium row the sharper signed constant instead gives
\begin{equation}\label{lt:medium-curvature}
 \Re K'(b)/H^2<(5/4)q^2+(23/2)q\mu,
\end{equation}
because $2/p_0^2<1250/24649$,
$6/5+(1250/24649)120000(9/556)^3<5/4$, and
$(1250/24649)225<23/2$.

\subsection{Centers retaining four moments}

In addition to $\theta$ put $\kappa=\sum\alpha_j^3$ and
$\lambda=-\sum\alpha_j^4$. Their bounds are
$|\kappa|\le q\mu$, $|\lambda|\le q\mu^2$. Assign weights
$1,2,3$ to $\theta,\kappa,\lambda$. For $s=\pm2\pi i$, the
weight-four center is
\begin{equation}\label{lt:four-center}
 Z_4=\zeta_s+u\kappa+v\theta\kappa+w\theta^2\kappa
                    +r\kappa^2+l\lambda+m\theta\lambda,
\end{equation}
where the four pure $\theta$ terms are those in
\eqref{lt:gaussian-center}, and
\begin{align*}
 u&=-s^3/3+s^2-2s,&
 v&=5s^4/6-13s^3/3+12s^2-20s,\\
 w&=-7s^5/4+157s^4/12-51s^3+131s^2-210s,&
 r&=s^5/3-13s^4/6+26s^3/3-22s^2+40s,\\
 l&=-s^4/4+s^3-3s^2+6s,&
 m&=3s^5/4-19s^4/4+18s^3-51s^2+90s.
\end{align*}
For higher weights define
\[
 J_4(z)=\int_0^ze^{w+\theta w^2/2+\kappa w^3/3+\lambda w^4/4}\,dw.
\]
Let $T_j(z)=\sum_{k=0}^j(-1)^{j-k}j!z^k/k!$. Its formal expansion is
\[
 J_4(z)=\sum_{a,b,c\ge0}
 \frac{\theta^a\kappa^b\lambda^c}{2^a3^b4^c a!b!c!}
 \{e^zT_{2a+3b+4c}(z)-T_{2a+3b+4c}(0)\}.
\]
Starting from $Z_0=s$, define recursively
\begin{equation}\label{lt:center-recursion}
 Z_j=Z_{j-1}-[\text{weight }j]J_4(Z_{j-1}),
        \qquad 1\le j\le8,\qquad e^s=1.
\end{equation}
Here $[\text{weight }j]$ selects the homogeneous component of weight $j$.
The new correction has derivative coefficient one, so this cancels
all coefficients up to the indicated weight. The first four levels
give exactly \eqref{lt:four-center}. The numbers of new terms in
levels $1,\ldots,8$ are $1,2,3,4,5,7,8,10$. Thus
\eqref{lt:center-recursion} uniquely specifies all forty nonconstant
terms of $Z_8$ as rational polynomials in $s$.

For $Z=Z_m(s)$, where $m=4$ or $8$ and $s=\pm2\pi i$, put
\begin{align*}
 A(t)&=t(Z-s)+\theta t^2 Z^2/2+\kappa t^3Z^3/3+\lambda t^4Z^4/4,\\
 A_r(t)&=(t-1)(Z-s)+\theta(t^2-1)Z^2/2
                  +\kappa(t^3-1)Z^3/3+\lambda(t^4-1)Z^4/4.
\end{align*}
These exponents give
\[
 J_4(Z)=Z\int_0^1e^{st+A(t)}\,dt,\qquad
 \frac{J_4(Z)}{J_4'(Z)}
 =Z\int_0^1e^{s(t-1)+A_r(t)}\,dt.
\]
For caps $q_0,\mu_0>0$ with $q\le q_0$, $\mu\le\mu_0$, use
the scaled variables
\[
 (u_1,u_2,u_3)=
 \left(\frac{\theta}{q_0},\frac{\kappa}{q_0\mu_0},
                         \frac{\lambda}{q_0\mu_0^2}\right),
 \qquad |u_1|,|u_2|,|u_3|\le1.
\]
For $\nu\in\mathbb Z_{\ge0}^3$, write
$\mathbf u^\nu=u_1^{\nu_1}u_2^{\nu_2}u_3^{\nu_3}$ and
$\mathrm{wt}(\nu)=\nu_1+2\nu_2+3\nu_3$; $\gamma\le\nu$ means
coordinatewise comparison. For either exponent $A$ or $A_r$, let
$A_\gamma(t)$ be its coefficient of $\mathbf u^\gamma$ after this
substitution, and let $E_\nu(t)$ be the coefficient of $\mathbf u^\nu$
in its exponential. Each $A_\gamma$ has degree at most four in $t$.
The exponential coefficients satisfy
\begin{equation}\label{lt:euler-recursion}
 E_0=1,\qquad
 E_\nu=\frac1{\mathrm{wt}(\nu)}
       \sum_{0<\gamma\le\nu}\mathrm{wt}(\gamma)A_\gamma E_{\nu-\gamma},
 \qquad \nu\ne0.
\end{equation}
The sum is over nonzero multi-indices $\gamma\le\nu$.
This identity is obtained by applying the weighted Euler derivative to
the chosen exponential. Since $e^s=1$, both integrals use the recurrence
for $I_m$ above. The integrated coefficients through weight $m$ vanish
exactly.
For the absolute residual $J_4(Z)$ or the relative residual
$J_4(Z)/J_4'(Z)$, let $F_*$ be the sum of the moduli of its
coefficients in $\mathbf u$ with $m<\mathrm{wt}(\nu)\le M$,
where $M=14$ for $m=4$ and $M=28$ for $m=8$.

On every interval
$[j/128,(j+1)/128]$ express each quartic coefficient in Bernstein
form. Its values lie in the convex hull of its five complex control
coefficients, so the largest control modulus bounds the entire
interval. The same applies to its difference from the endpoint value.
Alternatively the midpoint Taylor bound
$\sum_{k=0}^4|f^{(k)}(m_t)|h^k/k!$, where
$m_t=(2j+1)/256$ and $h=1/256$ are the midpoint and half-length,
bounds that interval. The four-moment rows use the former method and the
weight-eight rows use the latter. There are 253 path polynomials for
$Z_4$ and 1409 for $Z_8$.

For each exponent coefficient take the largest bound over the 128
intervals. Let $a_j$ be the sum of these bounds in weight $j$,
and $z_j$ the sum of the moduli of the coefficients of $Z$ in
$\mathbf u$ in that weight, including $z_0=|s|$. Put
$A_* (v)=\sum a_jv^j$ and $Z_*(v)=\sum z_jv^j$. The infinite tail
after weight $M$ is bounded by
\begin{equation}\label{lt:positive-weight-tail}
 T_*=Z_*(1)e^{A_*(1)}-
          \sum_{j=0}^M[v^j]\{Z_*(v)e^{A_*(v)}\}.
\end{equation}
The positive exponential coefficients are generated by
$e_0=1$, $e_j=j^{-1}\sum_{k=1}^jka_ke_{j-k}$.
All coefficients in the discarded tail are nonnegative. Thus
\eqref{lt:positive-weight-tail} bounds every omitted weight, including
every mixed parameter term. Its exponential is bounded upward by
\eqref{lt:exp-upper}, or by an outward complex ball evaluation with
an explicit remainder bound.

To pass from the model to the actual finite product, put
\[
 r_5=\frac{q_0\mu_0^3Z_*(1)^5}{5(1-\mu_0Z_*(1))},
       \qquad |Z|\le Z_*(1)<1/\mu_0.
\]
An absolute model residual $F_*+T_*$ has actual derivative-relative
bound $(F_*+T_*)e^{M_c+r_5}$, where $M_c$ bounds the center exponent
after $s$. A relative model residual instead has bound
$(F_*+T_*)e^{r_5}$. In both cases the actual product discrepancy is
bounded by
\begin{equation}\label{lt:four-actual-tail}
 \frac{q\mu^3Z_*(1)^6e^{M_r+2r_5}}{30(1-\mu_0Z_*(1))}.
\end{equation}
Here $M_r$ bounds the relative exponent along the path, and $1/30$
includes the integral of $t^5$. The path remainder and endpoint
denominator give the two $r_5$ terms.

There are two scaling rules. If $q=tq_0$, $0<t\le1$, and
$\mu\le\mu_0$, a monomial $\theta^a\kappa^b\lambda^c$ is bounded by
$t^{a+b+c}$ times its corner majorant. The uncancelled $Z_4$ terms
have total degree at least two; the tail above weight fourteen has
total degree at least five.
For $Z_8$ the corresponding degrees are at least three and ten.
These statements follow from $\mathrm{wt}(a,b,c)\le3(a+b+c)$
and the exact low-weight cancellation. They hold term by term in the
multivariate majorant and hence for its infinite tail. When instead
$q\le k/D$, $\mu=1/D$ and $D$ increases, the three parameters
scale with weights $1,2,3$.

\begin{proposition}\label{lt:q04mu04}\label{lt:q05mu02}\label{lt:W8}
For each row and $0<q\le q_0$, $0<\mu\le\mu_0$, the disks
$|S(z-b)-Z_m(s)|<R$, $s=\pm2\pi i$, are disjoint from each other
and from $b$, and contain exactly one simple original zero each,
denoted $z_\pm$ respectively.
Put $w_\pm=S(z_\pm-b)$. Both disks lie in the critical-free
disk $|S(z-b)|<Z_{\mathrm{out}}$. The final column gives
$|\Im S|<kQ/H$ at a rightmost zero.
\begin{center}\small
\begin{tabular}{ccclcc}
\toprule
$m$&$q_0$&$\mu_0$&$R(q,\mu)$&$Z_{\mathrm{out}}$&$k$\\\midrule
4&$1/20$&$1/50$&$9q$&$17/2$&$7$\\
4&$1/25$&$1/25$&$(45/4)q$&$17/2$&$15/2$\\
4&$3/70$&$1/70$&$34q^2+13400q\mu^3$&$17/2$&$21/4$\\
8&$3/50$&$1/50$&$64q^2+95000q\mu^3$&$19/2$&$8$\\
8&$273/5000$&$91/5000$&$22q^2+45000q\mu^3$&$9$&$33/5$\\\bottomrule
\end{tabular}
\end{center}
The first row also gives $|w_\pm|<17/2$ and
$|(w_++w_-)/2|<19/10$ for the same pair.
On the subrectangle $q\le57/1000$, $\mu\le1/51$, the fourth row gives
\begin{equation}\label{lt:Q057}
 |w_\pm|<8917/1000,\qquad
 |(w_++w_-)/2|<2109/1000,\qquad |\Im S|<(37/5)Q/H.
\end{equation}
\end{proposition}
\begin{proof}
Use \eqref{lt:center-recursion}--\eqref{lt:four-actual-tail}. For the
first two rows the exact caps of the center modulus are respectively
$8.030851273687$ and $7.706829806638$. Upper bounds at
the corners are as follows:
\begin{center}\small
\begin{tabular}{c|rrrrr}
\toprule
$(q_0,\mu_0)$&$F_*$&$T_*$&model/actual derivative&product error&circle error\\\midrule
$(.05,.02)$&$.091297$&$.017192$&$.216053$&$.008556$&$.200791$\\
$(.04,.04)$&$.079034$&$.012556$&$.169285$&$.049027$&$.200791$\\\bottomrule
\end{tabular}
\end{center}
The sums in the last three columns are strictly less than $9/20$.
The containing-disk bound is $L=31/20$ in both cases.
The finite sum retains all weights five through fourteen; the
infinite sum is exactly bounded by \eqref{lt:positive-weight-tail}.
The path, derivative and circle inequalities follow from
the deterministic recurrences and all 128 closed path intervals
specified above. At smaller $q$, both actual residuals are at most
$t=q/q_0$ times their corner values. Since
$(e^{LR}-1-LR)/L$ is convex and vanishes at zero, the circle error
is at most $t$ times its corner value. Thus the radii are $t(9/20)$
throughout the parameter rectangles.

For the third row the same construction at its own corner yields
\begin{gather*}
 Z_*(1)<7.629836,\quad |Z_4-s|<1.346650,\\
 M_c<.478584,\quad M_r<.481296,\quad
 F_*<.033415,\quad T_*<.000911.
\end{gather*}
Using the unrounded rational witnesses produced by the path conversion,
the model contribution is less than $30.181q^2$ and the actual
product contribution is less than $11961q\mu^3$.
Both are below nine tenths of the corresponding radius coefficient.
At the corner $R_0<.064124$, $Z_*(1)+R_0<17/2$,
\eqref{lt:log-derivative} is less than $31/20$, and the nonlinear
fraction is less than $.052<1/10$. The model terms and positive tail
are bounded by $(q/q_0)^2$ times their corner values;
\eqref{lt:four-actual-tail} retains its explicit $q\mu^3$.
Thus the third row holds throughout its parameter rectangle.

For the last two rows use the relative exponent $A_r$.
The upper bounds, including all weights at least twenty-nine,
are
\begin{center}\small
\begin{tabular}{c|rrrrrr}
\toprule
row&$F_*$&$T_*$&model error&product error&$L$&nonlinear fraction\\\midrule
$q_0=.06$&$.136221$&$.012814$&$.150033$&$.028511$&$7/4$&$1/3$\\
$q_0=.0546$&$.045190$&$.000674$&$.046017$&$.010525$&$17/10$&$1/5$\\\bottomrule
\end{tabular}
\end{center}
In each row the model error divided by $q_0^2$ is less than
$(1-\beta)$ times the $q^2$ radius coefficient, and the product
error divided by $q_0\mu_0^3$ is less than $(1-\beta)$ times the
other coefficient. The nonlinear fraction is less than
$\beta$. Containment, source exclusion and separation hold at both
corners and improve on the smaller rectangles. Lemma~\ref{lt:circle}
therefore applies.

The coefficient polynomials are generated by
\eqref{lt:center-recursion}; their exponential coefficients are given
by \eqref{lt:euler-recursion}; every path coefficient is bounded on
the stated 128 intervals; and the infinite positive tail is
\eqref{lt:positive-weight-tail}. The electronic appendix contains the
exact rational coefficient lists and programs performing these
operations at 352 or 448 bits with outward bounds. A proposed rational
modulus bound is accepted only if it exceeds its whole interval bound.
The vanishing identities are checked separately over $\mathbb Q[s]$.

Every center correction has positive total parameter degree, and
$R(q,\mu)\le(q/q_0)R_0$. If $\Gamma$ bounds $|Z_m-s|$ at
the corner, rightmostness gives
$2\pi|\Im\sigma|<(q/q_0)(\Gamma+R_0)$.
The rational or outward strict gates
$\Gamma+R_0<2\pi q_0k$ give exactly the last column of the table.
Averaging the two centers cancels all odd powers of $s$. The sum of
the absolute even coefficients at $(.05,.02)$, plus $9/20$, is
less than $1.890068<19/10$. At $(.057,1/51)$ the corresponding
evaluations give respectively
$8.916593$, $2.108213$ and $7.352979$, proving
\eqref{lt:Q057} for the pair $(z_+,z_-)$.
\end{proof}

\subsection{Curvature and the hole field}

\begin{lemma}\label{lt:hole-field}
Assume the unit-disk strict-hole geometry of
Lemma~\ref{pr:edge-reduction}. Suppose $C_0\ge0$ and
\[
 Qd/H<k,\qquad |\Im S|<K_0Q/H,\qquad
 \Re K'(b)<C_0/d^2.
\]
Then
\begin{equation}\label{lt:general-hole-field}
 H>\frac{N-(2C_0+2K_0k)/d}{2(1+a)}.
\end{equation}
\end{lemma}
\begin{proof}
Put $p_b=1-|b|^2\ge0$, $\lambda_b=1-|b-a|^2>0$.
For $v_j=(b-\xi_j)^{-1}$, the strict hole gives
\[
 \lambda_b|v_j|^2+2\Re((b-a)v_j)<1.
\]
This is an open disk whose greatest real coordinate is
$m_R=(1-x+a)/\lambda_b$. Rightmostness gives $\Re v_j\ge0$,
so
\[
 Q-2\sum(\Re v_j)^2=\Re K'(b)<C_0/d^2,\qquad
 Q<2m_RH+C_0/d^2.
\]
Since $\lambda_b\ge p_b$, one has $p_bm_R\le1-x+a$.
Also $p_b<2d$, because $d>1-|b|$ in the strict crescent.
Apply \eqref{pr:radial-identity} and $0\le y\le1$ to obtain
\[
 N<2(x+p_bm_R)H+2K_0Q/H+p_bC_0/d^2
   \le2(1+a)H+(2K_0k+2C_0)/d.
\]
This also covers $p_b=0$ without division by it.
\end{proof}

\begin{proposition}\label{lt:D70}
If $D\ge70$, then
\[
 |\Im S|<(21/4)Q/H,\qquad Qd/H<Y_{70}(D)<8/3,
\]
\[
 Y_{70}(D)=\frac{4D}{D-91/10+\sqrt{(D-91/10)^2-24D}}.
\]
In the strict-hole unit-disk geometry,
\begin{equation}\label{lt:field70}
 H>\frac{N-(596/5)/d}{2(1+a)}.
\end{equation}
If $D\ge100$, the stronger consequence is
\begin{equation}\label{lt:field100}
 H>\frac{N-(1899/20)/d}{2(1+a)}.
\end{equation}
\end{proposition}
\begin{proof}
First apply Proposition~\ref{mf:D25}, giving $Qd/H<3$.
This enters the third row of Proposition~\ref{lt:W8}.
The six additional moment terms give
$|Z_4-\zeta_s|<175q\mu$ on its parameter rectangle.
This follows by taking the modulus of each displayed coefficient in
\eqref{lt:four-center} at $s=2\pi i$; their total at the corner
is less than $.107054<175(3/70)(1/70)$.
The effective Gaussian radius is therefore
\[
 R_{\rm eff}=34q^2+(175+13400/4900)q\mu.
\]
Lemma~\ref{lt:signed-moment} with $6/5$ and the exact comparisons
\[
 6/5+(1250/24649)34<3,\qquad
 (1250/24649)(175+13400/4900)<91/10
\]
give $\Re K'(b)/H^2<3q^2+(91/10)q\mu$.
On the other hand $\sum(\Re v_j)^2\le H/d$, so
$\Re K'(b)/H^2\ge q-2\mu$.
With $y_0=Qd/H$ this yields
\[
 3y_0^2-(D-91/10)y_0+2D>0.
\]
At $D=70$ the polynomial is positive at $2$, negative at $3$,
decreasing on $[2,3]$, and equals $-16/15$ at $8/3$.
These signs persist for all $D\ge70$; the bound $y_0<3$ removes
the upper quadratic branch. The smaller root is $Y_{70}(D)$.
For $Y_{70}=2+v$ its inverse relation is
$D=151/(5v)+211/10+3v$, decreasing on $0<v\le2/3$.
This proves the assertion for all $D\ge70$.
For the hole field use $k=8/3$, $K_0=21/4$ and
$C_0=3k^2+(91/10)k=228/5$ in Lemma~\ref{lt:hole-field}.
The resulting constant is $596/5$. For $D\ge100$,
Proposition~\ref{mf:D100} gives $k=9/4$ instead.
Then $C_0=2853/80$ and the constant is $1899/20$.
\end{proof}

\begin{proposition}\label{lt:medium139}
If $D\ge139$, then
\[
 |\Im S|<4Q/H,\qquad
 Qd/H<Y(D):=\frac{4D}{D-23/2+\sqrt{(D-23/2)^2-10D}}<9/4.
\]
In the unit-disk strict-hole geometry,
\begin{equation}\label{lt:medium-field}
 H>\frac{N-(2637/32)/d}{2(1+a)}.
\end{equation}
\end{proposition}
\begin{proof}
Proposition~\ref{mf:D100} first gives $Qd/H<9/4$, which
enters the medium row of Proposition~\ref{lt:quartic}.
Equation~\eqref{lt:medium-curvature} gives
\[
 (5/4)y_0^2-(D-23/2)y_0+2D>0,\qquad y_0=Qd/H.
\]
At $D=139$ this is positive at $2$, negative at $9/4$, and
decreases on $[2,9/4]$; the signs persist on the degree-independent
ray. The bound $y_0<9/4$ excludes the upper branch. The inverse relation
at $Y=2+v$ is $D=28/v+33/2+(5/4)v$, decreasing for
$0<v\le1/4$. Finally apply Lemma~\ref{lt:hole-field} with
$k=9/4$, $K_0=4$ and $C_0=(5/4)k^2+(23/2)k$.
Its constant is $2637/32$.
\end{proof}

\begin{proposition}\label{lt:W8-field}
Suppose $Qd/H<3$ and $D\ge5000/91$. Then the second weight-eight
row gives $|\Im S|<(33/5)Q/H$ and
$\Re K'(b)<85/d^2$. In the strict-hole unit-disk geometry it follows
that
\begin{equation}\label{lt:field210}
 H>\frac{N-210/d}{2(1+a)}.
\end{equation}
\end{proposition}
\begin{proof}
The assumed ratio and field enter $q\le273/5000$, $\mu\le91/5000$.
The difference $Z_8-\zeta_s$ has coefficient-norm sum less
than $.397022$ at this corner. Adding the radius gives a
total effective Gaussian radius less than $.477420$.
Every term of this center difference has weight at least two.
As $D$ increases, with $t=(5000/91)/D$, its parameters lie in the
corner polydisk scaled by $t,t^2,t^3$. Consequently the effective
radius is at most $t^2$ times its corner bound.
Since $33/5<8$, the phase hypothesis for the $13/10$ bound in
Lemma~\ref{lt:signed-moment} holds. The coefficient sum in
\eqref{lt:center-recursion} gives
\[
 D^2\Re K'(b)/H^2<84.718<85.
\]
This comparison uses all forty terms and the whole radius; the
coefficient bound is a finite directed evaluation of those explicitly
defined polynomials. Lemma~\ref{lt:hole-field} with $k=3$,
$K_0=33/5$ and the sharper displayed coefficient gives a field
constant less than $2(84.718)+6(33/5)<210$.
\end{proof}

\section{Inverse logarithmic counting and transport}\label{sec:inverse-tools}

Use \eqref{pr:field-notation} and \eqref{pr:normalized-moments},
with $b$ an actual simple zero and $S\ne0$.

\subsection{Inverse rectangles}

\begin{lemma}\label{lt:inverse-log}\label{lt:odd-pi}
Suppose $Q/H<V$, $d\ge d_0>R>0$ and $H\ge H_0>0$. Put
\begin{gather*}
 \Lambda=\frac{VR}{1-R/d_0},\qquad
 e_R=V\left(\frac{R^2}{2}+\frac{R^3}{3(d_0-R)}\right),\qquad
 \widehat R=R-e_R,\\
 B_R=\frac{V}{(1-R/d_0)^2(1-\Lambda)^3},\qquad
 p_*=(1+\Lambda)^{-1},\qquad A_*=(1-\Lambda)^{-1}.
\end{gather*}
Assume $0<\Lambda<1$ and $\widehat R>0$. Let $0<\nu<1$, $W>0$,
and choose positive $e,\epsilon,g$ such that
\begin{gather}
 e>e^{-W},\qquad B_R/H_0<\epsilon,\qquad
 \nu B_R\widehat R<g,\qquad
 W^2<(1-\nu^2)(H_0\widehat R)^2,\label{lt:inverse-corner}\\
 \epsilon(1+e)+ge<1-A_*e,\quad
 \epsilon(1+e)+ge^2<p_*-e,\quad
 \epsilon+ge<p_*.\label{lt:inverse-sides}
\end{gather}
For any positive odd integer $\ell$ with
$\ell\pi\le\nu H\widehat R$, the disk $|z-b|<R$ contains at
least $\ell$ distinct simple original zeros. If
$\nu H\widehat R\ge\pi$, its number of original zeros is
strictly greater than
\begin{equation}\label{lt:inverse-density}
 \nu H\widehat R/\pi-2>\nu(7/22)H\widehat R-2.
\end{equation}
The source $b$ is included once in this count.
\end{lemma}
\begin{proof}
Let $L(z)=\log(p'(z)/p'(b))$ on $|z-b|<d$, with $L(b)=0$.
The finite logarithms give on $|z-b|\le R$
\[
 |L'(z)/S-1|<\Lambda,\qquad
 |L(z)-S(z-b)|<He_R.
\]
For the second bound retain the quadratic term and bound the entire
tail from order three by the displayed geometric series. The first
bound integrated along each joining segment proves injectivity: the
quotient $[L(z)-L(w)]/[S(z-w)]$ lies in $|u-1|<\Lambda$.
Rouch\'e's theorem on $|z-b|=R$ shows that the image contains
$|v|<H\widehat R$. Its analytic inverse $\phi$ satisfies
$\phi(0)=b$. Put $A(v)=S\phi'(v)$. Differentiation gives
\[
 A(0)=1,\qquad \Re A>p_*,\qquad |A|<A_*,\qquad |A'|<B_R/H.
\]
The last assertion uses
$|L''|\le Q/(1-R/d_0)^2$ and
$A'=-S L''(\phi)/(L'(\phi))^3$.

The whole rectangle $|\Re v|\le W$, $|\Im v|\le Y=\ell\pi$
lies strictly inside the inverse disk by \eqref{lt:inverse-corner}.
The actual primitive is
\[
 \frac{Sp(\phi(v))}{p'(b)}
  =\int_0^ve^uA(u)\,du=e^vA(v)-1-E(v),\qquad
 E(v)=\int_0^ve^uA'(u)\,du.
\]
Use the broken path $0,-W,-W+i\Im v,v$. For $v=x+iy$,
\[
 |E(v)|<\epsilon(1+e^x)+ge^{-W}.
\]
For $G=e^vA-1$, the left side has $|G|>1-A_*e^{-W}$,
the right side has $e^{-W}|G|>p_*-e^{-W}$, and each horizontal
side has $|G|>1+p_*e^x$, since $\ell$ is odd. Equations
\eqref{lt:inverse-sides} give $|E|<|G|$ on all four sides,
including every corner. Hence Rouch\'e compares the actual
primitive with $G$.

For the count, use the homotopy $A_t=1-t+tA$,
$0\le t\le1$. Then $\Re A_t\ge p_*$ and $|A_t|\le A_*$.
The same left and right estimates, and the strict negative real part
on the odd-$\pi$ sides, keep $e^vA_t-1$ nonzero throughout the
boundary. It has the same number of zeros as $e^v-1$, namely
$\ell$. Injectivity of $\phi$ and the nonvanishing of $p'$ in
the physical disk give distinct simple original zeros.
Finally, the largest permitted odd integer is
$1+2\lfloor(\nu H\widehat R/\pi-1)/2\rfloor$, which is
strictly larger than $\nu H\widehat R/\pi-2$.
\end{proof}

\begin{lemma}\label{lt:quarter-pi}
Use the notation and hypotheses of Lemma~\ref{lt:inverse-log},
replacing the third inequality in \eqref{lt:inverse-sides} by
\[
 \epsilon+ge<\frac{p_*}{1+p_*+A_*}.
\]
For an integer $k\ge0$ with $Y=(2k+1/2)\pi\le\nu H\widehat R$,
put $\mathcal R_Y=\{v:|\Re v|<W,\ |\Im v|<Y\}$.
Then $\phi(\mathcal R_Y)\subset\{z:|z-b|<R\}$ contains $2k+1$
simple original zeros. In particular, when
$\nu H\widehat R/\pi\ge1/2$, the number of original zeros in
$|z-b|<R$ is strictly greater than
$\nu H\widehat R/\pi-3/2$.
\end{lemma}
\begin{proof}
On these horizontal sides $e^{iy}=\pm i$. For $u=e^x>0$,
\[
 |e^vA(v)-1|\ge\max(1-A_*u,p_*u),\qquad
 \frac{\max(1-A_*u,p_*u)}{1+u}\ge\frac{p_*}{1+p_*+A_*}.
\]
The two quotients meet at $u=1/(p_*+A_*)$. This gives the
horizontal Rouch\'e comparison. On $y=\pm Y$, the amplitude homotopy
satisfies
\[
 \Im(e^vA_t-1)=e^x\sin y\,\Re A_t,
\]
which is positive on the upper side and negative on the lower side.
Here $\Re A_t>0$.
The model $e^v-1$ has exactly
$2k+1$ zeros. Choose the largest permitted $k$ to obtain the
last bound.
\end{proof}

The same inverse lemmas allow a sharper logarithmic radius. For any
integer $M\ge2$ one may replace $e_R$ by
\begin{equation}\label{lt:higher-log-loss}
 V\left(\sum_{k=2}^M\frac{R^k}{k d_0^{k-2}}
       +\frac{R^{M+1}}{(M+1)d_0^{M-1}(1-R/d_0)}\right).
\end{equation}
This retains every logarithmic order beyond $M$; all the other
constants and the full-boundary proof are unchanged. To use exact
integer counts, it is sufficient to verify
$\ell\pi_{\rm up}\le\nu H_0\widehat R$ for a rational upper
bound $\pi_{\rm up}>\pi$. If the source root is inserted separately
in a later weighted sum, use $\ell-1$ other roots, since $b$ was
already counted. The one-sided and narrow-strip refinements needed
near the origin are proved in Section~\ref{sec:origin}.

\subsection{Schur coupling and positive cumulative sums}

\begin{lemma}\label{lt:schur}
There is a nonnegative matrix $(c_{ij})$, indexed by all original
zeros and all critical points, with
\[
 \sum_jc_{ij}=1/n,\qquad \sum_ic_{ij}=1/N,\qquad
 N\sum_ic_{ij}z_i=\xi_j.
\]
If $|z_i|\le1$, then
$N\sum_ic_{ij}|z_i-\xi_j|^2\le1-|\xi_j|^2$.
\end{lemma}
\begin{proof}
Let $T=\diag(z_i)$, $e=n^{-1/2}(1,\ldots,1)^T$ and compress
$T$ to $e^\perp$. Its characteristic polynomial is $p'/n$ after
making $p$ monic. Indeed its cofactor in the $e$ direction is
$p(z)e^*(zI-T)^{-1}e=p'(z)/n$, first off the zeros and then
identically. Lift a unitary Schur basis of the compression to an
$n$ by $N$ matrix $U$. Its columns are orthonormal,
$UU^*=I-ee^*$, and the diagonal entries of $U^*TU$ are
$\xi_j$. Set $c_{ij}=|U_{ij}|^2/N$. All three identities follow.
Expanding around the barycenter gives
\[
 N\sum_ic_{ij}|z_i-\xi_j|^2
  =\sum_i|U_{ij}|^2|z_i|^2-|\xi_j|^2\le1-|\xi_j|^2.
\]
\end{proof}

\begin{lemma}\label{lt:transport}
Assume the unit-disk strict-hole geometry of
Lemma~\ref{pr:edge-reduction}, with the chosen simple rightmost zero
$b=x+iy$, $x\ge a$. For $0<a\le A<1$, define
\[
 \Gamma=\frac2{2-A^2},\quad u=\Gamma A,\quad
 q_A=\frac{u}{1-u},\quad E_A=\frac{2\Gamma}{(1-u)^2},
\]
and assume $u<1/2$. For $\xi_j=r_j\omega_j$, $|\omega_j|=1$, put
\[
 g=\frac1N\sum_j\frac1{b-\omega_j},\quad k=\Re g,\quad h=H/N,
 \quad \eta=\sum_{i,j}c_{ij}\frac{|z_i-\xi_j|^2}{|b-\xi_j|^2}.
\]
Then $k>0$, $h>(1-q_A)k$, and $\eta<E_Aak$.
\end{lemma}
\begin{proof}
The hole gives $1-a<r_j\le1$, and with
$t_j=a/2-\Re\omega_j$ it gives
\[
 (1-r_j)(1+r_j-a^2)<2ar_jt_j.
\]
In particular $t_j>0$. Since $2r/(1+r-a^2)$ increases on
$0<r\le1$, $1-r_j<\Gamma at_j$. Also
$t_j<x-\Re\omega_j\le|b-\omega_j|$. Hence
\[
 |b-r_j\omega_j|>(1-u)|b-\omega_j|,\qquad
 \left|\frac1{b-r_j\omega_j}-\frac1{b-\omega_j}\right|
 <q_A\Re\frac1{b-\omega_j}.
\]
Averaging gives $|S/N-g|<q_Ak$ and the field bound. The variance
part of Lemma~\ref{lt:schur}, together with
$1-r_j^2<2\Gamma at_j$, gives
\[
 \eta\le\frac1N\sum_j\frac{1-r_j^2}{|b-r_j\omega_j|^2}
 <\frac{2\Gamma a}{(1-u)^2N}
       \sum_j\frac{t_j}{|b-\omega_j|^2}\le E_Aak.
\]
\end{proof}

\begin{lemma}\label{lt:abel}
Suppose the number $L_j$ of original roots in nested disks
$|z-b|<t_jd$, where $0<t_1<\cdots<t_m<1$, satisfies
$L_j>s_jD-2$ with $0=s_0<s_1<\cdots<s_m$.
If $D\ge D_0>0$, $N\ge N_0$, and
\[
 c=\frac{N_0}{N_0+1}
 \left(\sum_{j=1}^m(1-t_j)^2(s_j-s_{j-1})-\frac2{D_0}\right)>0,
\]
then $\eta>chd$. Under Lemma~\ref{lt:transport},
$d<E_Aa/[(1-q_A)c]$.
\end{lemma}
\begin{proof}
Every root in the $j$th disk has
$|z_i-\xi_k|>(1-t_j)|b-\xi_k|$ for every $k$.
Its entire row in the coupling contributes more than
$w_j/n$, where $w_j=(1-t_j)^2$. Positive summation of cumulative
counts gives
\[
 n\eta>\sum_{j<m}(w_j-w_{j+1})L_j+w_mL_m
 >D\sum_jw_j(s_j-s_{j-1})-2w_1.
\]
Use $w_1<1$, $N/(N+1)\ge N_0/(N_0+1)$, and $D\ge D_0$.
Combining the result with $\eta<E_Aak<E_Aah/(1-q_A)$ proves
the last assertion. The cumulative counts enter the Abel sum with
nonnegative coefficients.
\end{proof}

For root-product weights, let $B_p=\prod_{i=1}^n f_i$, where
$0\le f_i\le1$ for every original zero, and set $-\log0=+\infty$.
If every root in the $j$th nested region has $f_i\le e^{-w_j}$,
where $w_1\ge\cdots\ge w_m\ge0$, and its cumulative count is
at least $K_j$, then
\begin{equation}\label{lt:abel-product}
 -\log B_p\ge w_mK_m+\sum_{j<m}(w_j-w_{j+1})K_j.
\end{equation}
This is the pointwise step-function lower estimate followed by
summation over every root. If $K_j=s_jH_*-2$, the budget is
$H_*\sum_jw_j(s_j-s_{j-1})-2w_1$.

\subsection{Crescent geometry and half-disk factors}

\begin{lemma}\label{lt:crescent}
Under the strict-hole geometry with $0<a\le x$, $b=x+iy$,
$y\ge0$, $|b|\le1$, put $t=\sqrt{1-a^2/4}$.
Every critical point satisfies
\begin{equation}\label{lt:crescent-gap}
 |b-\xi|^2>(x-a/2)^2+(t-y)^2.
\end{equation}
Consequently $d>x-a/2$, $d>1-y$, $d>1-|b|$, and
\begin{equation}\label{lt:quadrature-gap}
 d^2>(1-|b|)^2+a^2/4.
\end{equation}
\end{lemma}
\begin{proof}
Let $e=a/2+it$, write $\xi=e-u+iv$ and put $\lambda=y/t$.
Then $u>0$ by \eqref{pr:critical-halfplane}, and $0\le\lambda<1$
because $y\le\sqrt{1-x^2}\le\sqrt{1-a^2}<t$.
The exact identity
\[
 |b-\xi|^2-|b-e|^2
 =(1-\lambda)(|\xi-a|^2-1)
       +[2(x-a)+a\lambda]u+\lambda(u^2+v^2)>0
\]
proves the crescent comparison. Its difference from
$(1-y)^2$ is $x(x-a)+2y(1-t)\ge0$.
For fixed $r=|b|$ and $a$, the expression
$1+r^2-ax-2t\sqrt{r^2-x^2}$ increases for $a\le x<r$:
its derivative is $-a+2tx/\sqrt{r^2-x^2}>0$.
Continuity treats $x=r$. Thus its minimum occurs at $x=a$.
Since $2r-5a^2/4>0$ and
\[
 (2r-5a^2/4)^2-4(1-a^2/4)(r^2-a^2)
       =a^2(r-1)(r-4)+9a^4/16\ge0,
\]
squaring in the positive direction gives
\eqref{lt:quadrature-gap}. This includes $r=1$.
\end{proof}

\begin{lemma}\label{lt:direct-D70}
Under the preceding hypotheses put $c=x/2$. Then
\[
 D>\frac{Nc}{2\sqrt{2-2c+c^2}-2\sqrt2c(1-c)}.
\]
In particular, if $N\ge6250$ and $x\ge30501/500000$, then $D>70$.
\end{lemma}
\begin{proof}
For fixed $b$, the squared right side of \eqref{lt:crescent-gap}
decreases as $a$ increases to $x$, since its derivative is
$-x+ay/(2\sqrt{1-a^2/4})<0$ for $0<a\le x\le1$.
Put $t=\sqrt{1-c^2}$, $v=t-y\ge0$ and $G=\sqrt{c^2+v^2}$.
Thus $d>G$ and $p_b=1-|b|^2=-3c^2+2tv-v^2\le2v-3c^2-v^2$.
Since $\Re S>cQ$, the radial identity implies
\[
 D>\frac{NcG}{2|b|c+p_b}.
\]
The identities
$c+(1-c)v\le\sqrt{1+(1-c)^2}G$ and
$(3c^2+v^2)^2-8c^2G^2=(v^2-c^2)^2\ge0$,
together with $|b|\le1-p_b/2$, give
\begin{align*}
 2|b|c+p_b&\le2[c+(1-c)v]-(1-c)(3c^2+v^2)\\
 &\le
 [2\sqrt{2-2c+c^2}-2\sqrt2c(1-c)]G.
\end{align*}
The bracket is positive for $0<c\le1/2$.
For $c\ge c_0=30501/1000000$, one has
$2-2c+c^2<49/25<2$. Therefore the denominator is less than
$(14/5)(1-c+c^2)$ and
\[
 D>\frac{5N}{14}\frac{c}{1-c+c^2}
 \ge\frac{476578125000000}{6793005177007}>70.
\]
The rational function increases with $c$ because its derivative is
$(1-c^2)/(1-c+c^2)^2>0$.
\end{proof}

\begin{lemma}\label{lt:same-gap-field}
Assume all original zeros lie in $\overline\D$ and $a>0$ satisfies
$|a-\xi_j|>1$ for every critical point. Let $b=x+iy$ be a simple
rightmost zero with $x\ge a$, $y\ge0$, and $|b|\le1$, as in
Lemma~\ref{lt:crescent}. Put $c=x-a/2$ and $p_b=1-|b|^2$.
If $|\Im S|<K_0Q/H$ and $Nc>2K_0$, then
\[
 D>\frac{Nc^2-2K_0c}{2xc+p_b}.
\]
If $N\ge N_0$, $a\le a_U$, $x\ge x_0>a_U/2$, $p_b\le p_U$, and
$N_0(x_0-a_U/2)>2K_0$, then the uniform bound is obtained by
putting $N=N_0$, $c=x_0-a_U/2$, $x=x_0$, $p_b=p_U$.
\end{lemma}
\begin{proof}
The half-plane inequality gives $d>c$, $Q/H<1/c$. Hence
$N<(2x+p_b/c)H+2K_0/c$ and the first assertion follows.
With $x,p_b$ fixed, the derivative numerator in $c$ is
$2[Nxc^2+p_b(Nc-K_0)]>0$.
After writing $c=x-a_U/2$, the derivative numerator in $x$ is
$2[(Na_U/2+2K_0)c^2+p_b(Nc-K_0)]>0$.
The quotient decreases with $p_b$ and increases with $N$.
These signs also hold at $p_b=0$.
\end{proof}

\begin{lemma}\label{lt:half-disk-factor}
Let $b=x+iy$ be an actual simple rightmost zero, with
$x>0$, $y\ge0$, $0<B_L\le|b|\le B_U<1$,
$x\le X$, and $\gamma=b-a/2$, $0<a\le a_U$.
For a disk radius $R$ put
\[
 P=1-B_U^2,\qquad T=\frac{a_U}{2P},\qquad U=\frac RP,\qquad
 L=1+2TX.
\]
Suppose $R^2<P$ and $0<f<1$ satisfy
\begin{equation}\label{lt:half-disk-gate}
 f^2-(1-f^2B_L^2)(T^2+U^2)>2Uf^2LB_U.
\end{equation}
Then every original zero $z$ with $|z-b|\le R$ obeys
$|(\gamma-z)/(1-\overline z\gamma)|<f$.
\end{lemma}
\begin{proof}
Write $\tau=a/2$, $p_b=1-|b|^2$, $z=b+v$,
$L_\gamma=1-|\gamma|^2=p_b+2\tau x-\tau^2$ and
$A_0=p_b+\tau\overline b$. Rightmostness gives $\Re v\le0$.
The gate implies
\begin{gather*}
 T^2<\frac{f^2}{1-f^2B_L^2}<\frac1{1-B_L^2},\\
 \tau^2\le P^2T^2<\frac{P^2}{1-B_L^2}\le P\le p_b.
\end{gather*}
Thus $L_\gamma>0$, $|\gamma|<1$, and $1-f^2|\gamma|^2>0$.
Direct expansion gives
\[
 f^2|1-\overline z\gamma|^2-|\gamma-z|^2
 =f^2|A_0|^2-\tau^2-(1-f^2|\gamma|^2)|v|^2-2\Re(Vv),
\]
where $V=f^2L_\gamma\overline b+\tau(1-f^2)$.
Thus $\Re V>0$, $\Im V=-f^2L_\gamma y\le0$.
On the half disk $|v|\le R$, $\Re v\le0$,
$\Re(Vv)\le f^2L_\gamma yR$. The resulting lower bound is
\begin{align*}
 f^2p_b^2&-(1-f^2|b|^2)(\tau^2+R^2)
 +2f^2\tau x(p_b-R^2)+f^2\tau^2R^2-2Rf^2L_\gamma y.
\end{align*}
The two terms containing $2f^2\tau x$ and $f^2\tau^2$ are nonnegative.
After division by $p_b^2$, use $p_b\ge P>R^2$ and
\[
 \frac{L_\gamma}{p_b}
 =1+\frac{2\tau x-\tau^2}{p_b}\le1+2TX=L.
\]
The resulting bound is at least
$f^2-(1-f^2B_L^2)(T^2+U^2)-2Uf^2LB_U>0$.
The positive difference also excludes a zero denominator.
\end{proof}

For $N\ge N_0$, if a field estimate has the form $H>(N-c)/F$,
with $c\ge0$, $F>0$ and $(N_0-c)/F>H_*>0$, then
\begin{equation}\label{lt:energy-ray}
 \begin{gathered}
 \frac HN>\frac{1-c/N}{F}\ge\frac{1-c/N_0}{F}
             >\frac{H_*}{N_0},\\
 Q\ge\frac{H^2}{N}>\frac{N}{N_0^2}H_*^2\ge\frac{H_*^2}{N_0}.
 \end{gathered}
\end{equation}
The lower bound $(1-c/N)/F$ is nondecreasing in $N$, so these
estimates apply on the whole ray $N\ge N_0$.
